\chapter{拓扑向量空间.}

拓扑向量空间的一般理论大约创立于 1920 至 1930 年间。但在此之前，它早已由泛函分析中众多问题的研究作了长期的准备；而且，若不至少概括地指出这些问题的研究如何使数学家们（尤其从 20 世纪初起）逐渐认识到所论问题的共同渊源，以及把它们表述得远为一般、并对它们应用统一的求解程序的可能性，就无法追溯这段历史。

可以说，代数与分析之间的类比，以及把函数方程（也就是说未知量是一个函数的方程）看作代数方程的“极限情形”的想法，都追溯到无穷小演算的开端，这在某种方式上回应了这种“从有限到无限”的推广需要。但无穷小演算在代数上的直接祖先是有限差分演算（参见第 185 页以下），而不是一般线性方程组的求解；直到 18 世纪中叶，后者与微分学诸问题之间的最初类比才在弦振动方程的背景下显现出来。我们在这里不打算深入这个问题的历史细节；但我们必须指出两个基本思想的出现，它们在下文中还会不断出现，而且两者似乎都应归功于 D. 伯努利。第一个思想在于把弦的振动看作 n 个质点组成的系统的振动当 n 趋于无穷时的“极限情形”；众所周知，当 n 有限时，这个问题在稍后给出了寻求线性变换特征值的第一个例子（参见第 88 页）；在所考虑的“过渡到极限”中，与这些数相对应的是弦的“固有振动”的频率，它们早在此前很久就已被实验观察到，而其理论上的存在性已在本世纪初（特别是由泰勒）建立。这个形式上的类比，尽管在后来的文献中相当少被提及（{[}302 b{]}，第 390 页），在 19 世纪期间似乎从未被忽视；但正如我们将在后文看到的，直到大约 1890–1900 年它才获得其全部重要性。

D. 伯努利的另一个想法（也许是受实验事实的启发）是“叠加原理”，按照这个原理，弦的最一般的振动必须能够“分解”为“固有振动”的叠加，用数学的语言说，这意味着弦振动方程的一般解必须能够展开为级数 \(\sum_{n} c_{n} \phi_{n}(x, t)\)，其中 \(\phi_{n}(x, t)\) 表示各个固有振动。众所周知，这一原理后来引发了一场关于把“任意”函数展开为三角级数之可能性的长期论战，这场论战直到 19 世纪前三分之一的时间里才由傅里叶和狄利克雷的工作了结。但甚至在这一结果达到之前，人们已经遇到过“正交”\(^{1}\)函数展开为级数的其他例子：球面函数和勒让德多项式，以及形如 \((e^{i\lambda_{n}x})\) 的各种系统，其中 \(\lambda_{n}\) 不再是同一个数的倍数，这些系统从 18 世纪起就被引入振动问题之中，也被傅里叶和泊松在关于热理论的研究中引入。大约在 1830 年，斯图姆{[}302 a and b{]}和刘维尔{[}204 a and b{]}把这些各别特殊情形中观察到的全部现象系统化为一个关于单变量函数的一般振动理论：他们考虑微分方程

\[
\frac {d}{d x} \left(p (x) \frac {d y}{d x}\right) + \lambda \rho (x) y = 0 (p (x) > 0, \rho (x) > 0)\tag{21.1}
\]

连同边界条件

\[
\left\{ \begin{array}{l l} y ^ {\prime} (a) - h _ {1} y (a) & = 0 \\ y ^ {\prime} (b) + h _ {2} y (b) & = 0 \end{array} \right. (h _ {1} \geq 0, h _ {2} \geq 0, a <   b)\tag{21.2}
\]

并证明以下基本结果：

\begin{enumerate}
\def\labelenumi{\arabic{enumi})}
\item
  这个问题具有 \(\neq 0\) 的解，仅当 \(\lambda\) 取到某个数列 \((\lambda_{n})\)（其中的数 \textgreater{} 0 且趋于 \(+\infty\)）中的值；
\item
  对每个 \(\lambda_{n}\)，解都是同一个函数 \(v_{n}\) 的倍数，这个函数可以假定已由条件 \(\int_{a}^{b} \rho v_{n}^{2} dx = 1\) 加以“规范化”，而且我们有 \(\int_{a}^{b} \rho v_{m} v_{n} dx = 0\)（对 \(m \neq n\)）；
\item
  每个函数 \(f\)，若它在 \([a, b]\) 中二次可微且满足边界条件 (21.2)，则可展开为一致收敛的级数 \(f(x) = \sum_{n} c_{n} v_{n}(x)\)，其中 \(c_{n} = \int_{a}^{b} \rho f v_{n} dx\)；
\item
  我们有等式 \(\int_{a}^{b}\rho f^{2}dx=\sum_{n}c_{n}^{2}\)（帕塞瓦尔已于 1799 年——而且是以纯粹形式的方式——对三角函数系证明过它，由它立即可得出后者（仍就三角级数）于 1828 年陈述的“贝塞尔不等式”）。
\end{enumerate}

半个世纪之后，这些性质由格拉姆的工作{[}133{]}加以补全；格拉姆继切比雪夫的研究之后，揭示了正交函数级数展开与“最佳平方逼近”问题（它直接源出于高斯在误差理论中的“最小二乘法”）之间的关系：给定函数的一个有限序列 \((\psi_{i})_{1\leq i\leq n}\)，这个问题的内容是，对函数 f 求线性组合 \(\sum_{i}a_{i}\psi_{i}\)，使积分 \(\int_{a}^{b}\rho(f-\sum_{i}a_{i}\psi_{i})^{2}dx\) 达到最小。原则上这只是一个基本的线性代数问题，但格拉姆以独创的方式解决了它，即对 \(\psi_{i}\) 施行通常以埃哈德·施密特之名著称的“规范正交化”程序。接下来转到无穷正交系 \((\phi_{n})\) 的情形，他给自己提出的问题是：用序列前 n 个函数的线性组合对函数 f 所作的“最佳平方逼近” \(\mu_{n}\)，何时当 n 趋于无穷时趋于 0；\(^{2}\)他由此被引导去定义完备规范正交系的概念，并认识到这一性质等价于不存在 \(\neq0\) 的、与所有 \(\phi_{n}\) 都正交的函数。他甚至试图阐明“均方收敛”的概念，但是，在测度论的基本概念被引入之前，他在这个方向上很难得到任何十分明确的结果。

在 19 世纪下半叶，分析学家们的主要力量主要用于把斯图姆–刘维尔理论推广到多变量函数，数理物理中椭圆型偏微分方程的研究以及自然伴随它们的边值问题尤其引导了这一推广。兴趣主要集中在“振动膜”方程

\[
L _ {\lambda} (u) \equiv \Delta u + \lambda u = 0\tag{21.3}
\]

其中所寻求的是在相当正规的区域 \(G\) 的边界上为零的解；这个问题所呈现的巨大分析困难只是一点一点地被克服的，人们原本认为不可能把对单变量函数行之有效的方法应用于这个问题。让我们回顾通向解决的主要阶段：\(G\) 的“格林函数”的引入，其存在性由施瓦茨证明；最小特征值的存在性的证明，同样出自施瓦茨；最后，1894 年，在一篇著名的论文中（{[}251 a{]}，第 IX 卷，第 123–196 页），H. 庞加莱成功地证明了全部特征值的存在性及其本质性质，其办法是：对给定的“右端” \(f\)，考虑这样的解 \(u_{\lambda}\)：它满足方程 \(L_{\lambda}(u) = f\) 且在边界上为零，并通过巧妙地推广施瓦茨的方法证明 \(u_{\lambda}\) 是复变量 \(\lambda\) 的亚纯函数，只有实的单极点 \(\lambda_n\)，它们正是所寻求的特征值。

这些研究与线性积分方程论的开端紧密相连，后者对现代观念之到来所作的贡献无疑最大。我们在这里只限于对这门理论的发展给出一些简短的说明。这类函数方程最初在 19 世纪上半叶零星地出现（阿贝尔、刘维尔），自从比尔和 C. 诺伊曼把相当正规的区域 G 上“狄利克雷问题”的求解化为一个“第二类积分方程”

\[
u (x) + \int_ {a} ^ {b} K (x, y) u (y) d y = f (x)\tag{21.4}
\]

的求解（未知函数为 u）；这个方程 C. 诺伊曼已于 1877 年借助“逐次逼近”程序解出。无疑既受到已提到的代数类比的驱使，也受到他刚刚在振动膜方程上所得结果的驱使，H. 庞加莱在 1896 年（{[}251 a{]}，第 IX 卷，第 202–272 页）想到在前述方程的积分之前引入一个变参数 \(\lambda\)，并断言，如同在振动膜方程的情形一样，解这时是 \(\lambda\) 的亚纯函数；但他没有能证明这一结果，它（对连续的“核” K 和有限区间 \([a, b]\)）直到四年后才由 I. 弗雷德霍姆建立{[}116{]}。后者也许甚至比他的前辈们更为自觉，让自己完全受 (21.4) 与线性方程组

\[
\sum_ {q = 1} ^ {n} \left(\delta_ {p q} + \frac {1}{n} a _ {p q}\right) x _ {q} = b _ {p} (1 \leq p \leq n)\tag{21.5}
\]

的类比的引导，以便把 (21.4) 的解表示成两个表达式之商，这两个表达式是仿照克拉默公式中出现的行列式构造出来的。这也不是一个新想法：早在 19 世纪初，“待定系数法”（其内容是：对一个假定可展开为级数 \(\sum_{n} c_{n} \phi_{n}\) 的未知函数——其中 \(\phi_{n}\) 是已知函数——通过计算系数 \(c_{n}\) 来求得它）就已引出了“无穷多个未知数的线性方程组”

\[
\sum_ {j = 1} ^ {\infty} a _ {i j} x _ {j} = b _ {i} (i = 1, 2, \dots).\tag{21.6}
\]

傅里叶碰到这样一个方程组时，仍像 18 世纪的数学家那样去“解”它：他把指标 i 或 j 大于 n 的项统统删去，用克拉默法则显式地解出这样得到的有限方程组，然后“过渡到极限”，在解中令 n 趋于 +∞！到后来，当这种偷懒的把戏不再够用时，对这一问题的最初尝试仍然是通过行列式理论作出的；从 1886 年起（继希尔的工作之后），H. 庞加莱、然后是 H. 冯·科赫，建立了一种“无穷行列式”理论，使某些类型的方程组 (21.6) 可以按古典的模式求解；即使这些结果不能直接适用于弗雷德霍姆所瞄准的问题，至少可以肯定，冯·科赫的理论特别曾作为构造他的“行列式”的范本。

正是在这个时刻希尔伯特登场，给这门理论以新的推动{[}163 b{]}。他首先完成弗雷德霍姆的工作，有效地实现了从 (21.5) 的解到 (21.4) 的解的过渡到极限；但他随即为实二次型的理论补上了相应的过渡到极限，具有对称核（也就是说 \(K(y, x) = K(x, y)\)）的那类积分方程自然导向这个理论，而且这类方程在数理物理中最为常见。他就这样得到了那个直接推广二次型约化到主轴的基本公式

\[
\int_ {a} ^ {b} \int_ {a} ^ {b} K (s, t) x (s) x (t) d s d t = \sum_ {n = 1} ^ {\infty} \frac {1}{\lambda_ {n}} \left(\int_ {a} ^ {b} \phi_ {n} (s) x (s) d s\right) ^ {2},\tag{21.7}
\]

其中 \(\lambda_{n}\) 是核 \(K\) 的（必然为实的）特征值，\(\phi_{n}\) 构成相应本征函数的正交系，而且当 \(\int_{a}^{b} x^{2}(s) dx \leq 1\) 时公式 (21.7) 的右端是一个收敛级数。他还表明，每个可“表示”为形式 \(f(x) = \int_{a}^{b} K(x, y) g(y) dy\) 的函数如何容许“展开” \(\sum_{n=1}^{\infty} \phi_{n} \int_{a}^{b} \phi_{n}(y) f(y) dy\)，而且循着与二次型经典理论的类比，他指出了一个用变分方法确定 \(\lambda_{n}\) 的程序，它无非是二次曲面主轴那个著名的极值性质的推广（{[}163 b{]}，第 1–38 页）。

希尔伯特的这些最初结果几乎立即被 E. 施密特以更简单也更一般的形式接手，避免了“弗雷德霍姆行列式”的引入以及从有限到无限的过渡，而且已经非常接近抽象的陈述，积分的线性性与正性这两个基本性质显然只在证明中被使用{[}274 a{]}。但希尔伯特那时已经达到了还要更为一般得多的观念。此前的一切工作凸显了平方可积函数的重要性，而帕塞瓦尔公式在这些函数与序列 \((c_{n})\)（它们满足 \(\sum_{n} c_{n}^{2} < +\infty\)）之间建立了紧密的联系。无疑正是这个想法引导着希尔伯特在 1906 年的论文（{[}163 b{]}，第 XI–XIII 章）中重新拾起古老的“待定系数法”，证明积分方程 (21.4) 的求解等价于一个由无穷多个线性方程组成的方程组

\[
x _ {p} + \sum_ {q = 1} ^ {\infty} k _ {p q} x _ {q} = b _ {p} (p = 1, 2, \dots)\tag{21.8}
\]

其中未知函数 u 的“傅里叶系数”为 \(x_{p} = \int_{a}^{b} u(t) \omega_{p}(t) dt\)（相对于一个给定的完备规范正交系 \((\omega_{n})\)）（其中 \(b_{p} = \int_{a}^{b} f(t) \omega_{p}(t) dt\)，\(k_{pq} = \int_{a}^{b} \int_{a}^{b} K(s, t) \omega_{p}(s) \omega_{q}(t) ds dt\)。而且，从这个观点看，(21.8) 中唯一值得考虑的解是那些使 \(\sum_{n} x_{n}^{2} < +\infty\) 的解；希尔伯特也系统地把自身局限于这种类型的解；但与此相对，他放宽了加于“无穷矩阵”（\(k_{pq}\)）上的条件（在 (21.8) 中该矩阵满足 \(\sum_{p,q} k_{pq}^{2} < +\infty\)）。从那时起，事情很清楚：由实数序列 \(x = (x_{n})\)（满足 \(\sum_{n} x_{n}^{2} < +\infty\)）组成的“希尔伯特空间”虽然未被明确引入，却潜藏于整个理论之下，并且显示为从有限维欧氏空间出发的一个“过渡到极限”。除此以外，对后继的发展特别重要的是，希尔伯特被引导在这个空间中引入的不仅是一种、而是两种不同的收敛概念（相应于所谓弱拓扑和强拓扑\(^{3}\)），以及一个“选择原理”，它无非是单位球的弱紧性。他在求解方程组 (21.8) 的背景下发展的新线性代数，完全建立在这些拓扑概念之上：线性映射、线性型和双线性型（与线性映射相关联）按照它们的“连续性”\(^{4}\)性质被分类和研究。特别地，希尔伯特发现弗雷德霍姆方法的成功依赖于“全连续”概念，他通过就双线性型对它进行表述\(^{5}\)并作深入研究来揭示这一概念；我们在这里无法就这一重要概念给出更多细节，也无法详述希尔伯特在其间开创对称双线性型（有界或否）谱论的那令人赞叹的深刻工作。

希尔伯特的语言仍是古典的，而且在整个《线性积分方程一般理论》中，他始终注视着理论的应用，并发展了众多的例子（约占全卷一半）。下一代人已经采取了更为抽象的观点。在弗雷歇和 F. 里斯关于一般拓扑的思想影响下（见第 142 页），E. 施密特{[}274 b{]}和弗雷歇本人在 1907–1908 年间有意把欧氏几何的语言引入“希尔伯特空间”（实的或复的）；正是在这些著作中第一次提到范数（用现在的记号 \(||x||\)）、它所满足的三角不等式、希尔伯特空间“可分”且完备的事实；此外，E. 施密特证明了到闭线性簇上的正交投影的存在性，这使他能把希尔伯特的线性方程组理论表述为更简单也更一般的形式。同样在 1907 年，弗雷歇和 F. 里斯注意到平方可和序列的空间有一个完全类似的“几何”；几个月后，当 F. 里斯和 E. 菲舍尔证明这个空间完备且同构于“希尔伯特空间”时，这一类比得到了完美的解释，同时也以令人惊异的方式凸显了勒贝格新创造的工具的价值。从那时起，希尔伯特空间理论的要点即可视为已被获得；在较晚近的进展中，最值得一提的是 M. H. 斯通和 J. 冯·诺伊曼大约在 1930 年给出的理论的公理化表述，以及大约 1934 年在雷利希、勒维希和 F. 里斯的工作中对“可分性”限制的摆脱{[}260 g{]}。

然而，在 20 世纪的最初几年，另外一些思想潮流也在涌现，以加强导向赋范空间理论的趋势。“泛函”的一般观念（也就是定义在一个集合上的取数值的函数，该集合的元素本身是一个或多个实变量的数值函数）出现于 19 世纪最后几十年，一方面与变分法相联系，另一方面与积分方程论相联系。如果说这一观念以及更一般的“算子”观念的揭示主要应归功于以平凯勒、尤其是沃尔泰拉为中心的意大利学派，那么这个学派的工作往往仍带有相当形式的性质，并且局限于特殊类型的问题，缺乏对底层拓扑概念足够深入的分析。1903 年，阿达马开创了“拓扑”对偶的现代理论，他寻求紧区间 I 上数值函数的空间 \(\mathcal{C}(I)\)（配备一致收敛拓扑的空间）上最一般的连续线性“泛函”，并把它们刻画为积分序列 \(x \mapsto \int_{I} k_{n}(t)x(t)dt\) 的极限。1907 年，弗雷歇和 F. 里斯以同样的方式表明，在希尔伯特空间上，连续线性型就是希尔伯特引入的“有界”型；随后在 1909 年，F. 里斯把阿达马定理定型，用斯蒂尔杰斯积分表达了 \(\mathcal{C}(I)\) 上的一切连续线性泛函，这个定理后来成为现代积分论的出发点（见第 226 页）。

第二年，还是 F. 里斯{[}260 c{]}在理论上作出了新的重要进展，他引入并研究了（仿照希尔伯特空间理论）区间 I 上 p 次可和函数的空间 \(L^{p}(I)\)（指数 p 满足 \(1 < p < +\infty\)），三年后他继之以{[}260 e{]}关于序列空间 \(L^{p}(\mathbf{N})\) 的类似工作；正如我们将在后文看到的，这些研究大大有助于澄清关于对偶的观念，因为在这里第一次遇到了两个互相对偶却并非自然同构的空间。\(^{6}\)

从那时起，F. 里斯就在考虑一个涵盖所有这些结果的公理化研究（{[}260 c{]}，第 452 页）；而且似乎只是出于一位分析学家不愿过分背离古典数学的顾虑，才使他没有以这种形式写出他 1918 年那篇关于弗雷德霍姆理论的著名论文{[}260 f{]}。他在那里原则上考虑紧区间上连续函数的空间 \(\mathcal{C}(I)\)；但在定义了这个空间的范数、并注意到 \(\mathcal{C}(I)\) 连同这个范数是完备的之后，他在论证中就再没有使用过完备赋范空间公理以外的任何东西。\(^{7}\)在这里不对这项工作作详细的考察，只须指出：正是在那里第一次以一般的方式定义了全连续线性映射的概念（通过把一个邻域变为相对紧集的性质）；\(^{8}\)而且，凭借公理化分析的一项杰作，弗雷德霍姆理论的全部（就其定性方面而言）被化归为单一的基本定理，即每个局部紧的赋范空间都是有限维的。

赋范空间的一般定义是 1920–22 年间由 S. 巴拿赫、H. 哈恩和 E. 黑利给出的（后者只考虑实数或复数序列的空间）。在其后的十年里，空间理论主要围绕两个对应用具有根本重要性的问题发展：对偶理论以及与贝尔“纲”概念相联系的诸定理。

我们已经看到，（拓扑意义上的）对偶观念可追溯到 20 世纪初；它潜存于希尔伯特的理论之下，并在 F. 里斯的工作中占据中心位置。例如后者早在 1911 年（{[}260 d{]}，第 41–42 页）就注意到，关系 \(|f(x)| \leq M||x||\)（在希尔伯特空间中被取作“有界”线性泛函的定义）当我们在空间 \(\mathcal{C}(I)\) 中时等价于 f 的连续性，而且所用的是一个具有完全一般性质的论证。在关于 \(\mathcal{C}(I)\) 上连续线性泛函的刻画方面，他还注意到，集合 A 在 \(\mathcal{C}(I)\) 中稠密的条件是：不存在 I 上与 A 的每个函数都“正交”的斯蒂尔杰斯测度 \(\mu \neq 0\)（从而推广了完备规范正交系的格拉姆条件）；最后，在同一项工作中他指出，空间 \(L^{\infty}\) 的对偶比斯蒂尔杰斯测度的空间“更大”（{[}260 d{]}，第 62 页）。

另一方面，在他关于空间 \(L^{p}(I)\) 和 \(L^{p}(\mathbf{N})\) 的工作中，F. 里斯成功地对 E. 施密特{[}274 b{]}给出的求解希尔伯特空间中线性方程组的方法作了修改，使之适用于更一般的情形。E. 施密特的想法在于，通过在由方程组 (21.6) 所表示的闭线性簇中寻求到原点的距离最小的点，来确定 (21.6) 的一个“极值”解。利用同一想法，F. 里斯证明，存在函数 \(x \in L^{p}(a, b)\) 满足方程组

\[
\int_ {a} ^ {b} \alpha_ {i} (t) x (t) d t = b _ {i} (i = 1, 2, \dots)\tag{21.9}
\]

（其中 \(\alpha_{i}\) 属于 \(L^{q}\)（满足 \(\frac{1}{p} +\frac{1}{q} = 1\)），而且还满足 \(\int_{a}^{b}|x(t)|^{p}dt\leq M^{p}\)），其必要充分条件是：对实数的所有有限序列 \((\lambda_i)_{1\leq i\leq n}\)，有

\[
\left| \sum_ {i = 1} ^ {n} \lambda_ {i} b _ {i} \right| \leq M. \left(\int_ {a} ^ {b} \sum_ {i = 1} ^ {n} | \lambda_ {i} \alpha_ {i} (t) | ^ {q} d t\right) ^ {1 / q}.\tag{21.10}
\]

1911 年{[}260 d{]}，他以类似的方式处理“广义矩问题”，即求解方程组

\[
\int_ {a} ^ {b} \alpha_ {i} (t) d \xi (t) = b _ {i} (i = 1, 2, \dots)\tag{21.11}
\]

其中 \(\alpha_{i}\) 连续而未知量是一个斯蒂尔杰斯测度 \(\xi;^{9}\)；这里可以清楚地看到，这个问题可以这样来陈述：按一个连续线性泛函在空间 \(\mathcal{C}(I)\) 中一列给定点上的值来确定这个泛函。黑利在 1912 年处理的也是这个形式的问题——他以一种颇为不同、适用范围更广的方法得到了 F. 里斯的条件\(^{10}\)——并在 1921 年以远为一般的条件重新处理了它。如我们上文所见，他引入了（序列空间上的）范数概念，并注意到这个概念推广了 n 维空间中凸体的“度规”，它被闵可夫斯基用于他关于“数的几何”的著名著作{[}221 a and b{]}。在这项工作中，闵可夫斯基还（在 \(R^{n}\) 中）定义了支撑超平面和“支撑函数”的概念{[}221 b{]}，并证明了在凸体的所有边界点处支撑超平面的存在性（{[}221 a{]}，第 33–35 页）。黑利把这些概念推广到配备任意范数的一个序列空间 E；他在 E 与空间 \(E'\)（由满足下列条件的序列 \(u = (u_n)\) 组成）之间建立了一种对偶：对每个 \(x = (x_n) \in E\)，级数 \((u_n x_n)\) 应当收敛；以 \(\langle u, x \rangle\) 记这个级数的和，他在 \(E'\) 中用公式 \(\sup_{x \neq 0} |\langle u, x \rangle| / ||x||\) 定义了一个范数，它在有限维空间中给出的正是支撑函数。\(^{11}\)E 中方程组 (21.6) 的解——其中每个序列 \(u_i = (a_{ij})_{j \geq 1}\) 都假定属于 \(E'\)——如黑利当时所证明的，化归为相继求解下列两个问题：\(1^\circ\) 在赋范空间 \(E'\) 上找一个连续线性型 L，使对每个指标 i 有 \(L(u_i) = b_i\)，如他所指出的，这导致 (21.10) 型的条件；\(2^\circ\) 寻求这样一个线性型能否写成 \(u \mapsto \langle u, x \rangle\)（对某个 \(x \in E\)）。正如黑利所观察到的，最后这个问题即使当 L 存在时也未必有解，他只限于给出一些在某些特殊情形蕴含解 \(x \in E\) 存在的充分条件{[}156{]}。

这些观念于 1927 年在 H. 哈恩的一篇基本论文{[}142{]}中获得了它们最终的形式，其结果两年后被 S. 巴拿赫（独立地）重新得到{[}14 c{]}。哈恩把闵可夫斯基–黑利程序应用于任意的赋范空间，从而赋予对偶空间一个（完备）赋范空间的结构；这使哈恩能立即考虑赋范空间的逐次对偶，并以一般的方式陈述黑利已预见到的自反空间问题。但尤其是，在保持范数的条件下延拓连续线性泛函这个压轴问题，由哈恩以完全一般的方式最终解决，用的是对维数作超限归纳的论证——由此给出了选择公理应用于泛函分析的一个最早的重要例子。\(^{12}\)在这些结果之外，巴拿赫又对连续线性映射与其转置之间的关系作了透彻的研究，借助一个关于对偶空间中弱闭子集的深刻定理，把直到那时只对空间 \(L^{p}\){[}260 c{]}已知的结果推广到一般的赋范空间；这些结果若使用几年后由豪斯多夫和巴拿赫本人引入的赋范空间的商空间概念来表述，就更为醒目。最后，又是巴拿赫发现了单位球的弱紧性（如我们上文所指出的，在许多特殊情形都被观察到）与自反性之间的联系，至少对可数型的空间是如此（{[}14 a{]}，第 189 页）。从此刻起，赋范空间的对偶理论就可以认为在其主要方面已经定型。

与此同时，一些外表悖谬的定理得到了澄清，其最早的例子可追溯到 1910 年前后。事实上，黑林格和托普利茨在那一年实质上证明了：希尔伯特空间上有界双线性型的一个序列 \(B_n(x,y)\)，若其值 \(B_n(a,b)\) 对每个给定的对 \((a,b)\) 都有界（界数先验地依赖于 \(a\) 和 \(b\)），则它在每个球上事实上都一致有界。他们的证明用反证法，其内容是用一种后来以“滑动驼峰法”著称的归纳法构造出一个违反假设的特殊对 \((a,b)\)，这个方法至今仍用于许多类似的问题。早在 1905 年，勒贝格就曾用类似的程序证明存在其傅里叶级数在某些点发散的连续函数；而且，在与黑林格和托普利茨同一年，他用同一方法证明了 \(L^1\) 中弱收敛的级数按范数有界。\(^{13}\)这类例子在随后的年代里成倍增加，但直到 1927 年都没有引入任何新观念；这一年，巴拿赫和施坦豪斯（有 S. 萨克斯的部分合作）把这些现象与贫集的概念以及完备度量空间上的贝尔定理联系起来，得到了一个涵盖此前所有特殊情形的一般命题{[}15{]}。此外，对完备赋范空间中“纲”问题的研究还同时把巴拿赫引向关于连续线性映射的许多其他结果；最引人注目也无疑最深刻的是“闭图象”定理，它与巴拿赫–施坦豪斯定理一样，在现代泛函分析中显示出是一阶工具{[}14 c{]}。

巴拿赫关于《线性算子》的论著的出版{[}14 a{]}，可以说标志着赋范空间理论成年期的开始。我们一直在谈论的所有结果以及许多其他结果都铺陈于这部书中，方式仍有些杂乱，但伴有取自分析各个领域的许多醒目的例子，它们似乎预示着这门理论辉煌的前景。事实上，这部著作取得了巨大的成功，其最直接的效果之一就是巴拿赫所用的语言和记号几乎被普遍采用。但是，尽管四十年来对巴拿赫空间进行了大量的研究（{[}203 bis{]}），如果不把巴拿赫代数理论及其在交换和非交换调和分析中的应用算作例外，那么这门理论对古典分析重大问题的几乎完全缺失的新应用，多少使人们对它寄予的希望落空了。

最为富有成效的发展，毋宁说是沿着对与赋范空间有关的概念作一种日益扩张、日益深入的公理化分析的方向发生的。虽然自 20 世纪初以来遇到的功能空间大都天生带有“自然的”范数，但也不是没有引起人们注意到一些例外。大约在 1910 年，E. H. 摩尔曾提议推广一致收敛的概念，即以“相对一致收敛”的概念来代替它，其中 0 的一个邻域由满足关系 \(|f(t)| \leq \epsilon g(t)\) 的函数 f 组成，g 是一个处处 \textgreater{} 0 的函数，它可以随邻域而变化。另一方面，在 1930 年之前人们已经注意到，像简单收敛、可测函数的依测度收敛或整函数的紧收敛这样的概念，都不能借助范数来定义；而且在 1926 年，弗雷歇已注意到这种性质的向量空间可以是可度量化的和完备的。但这些更一般空间的理论，只有与凸性思想相结合才能富有成效地发展。后者（我们在黑利那里已看到它的出现）曾是巴拿赫及其学生们研究的对象，他们认识到有可能以此把赋范空间理论中的许多命题解释得更具几何意味，从而为 J. 冯·诺伊曼 1935 年给出的局部凸空间的一般定义准备了道路。这些空间的理论，特别是涉及对偶的那些问题，最近十年间得到了特别的发展。在这方面必须指出的，一方面是一般拓扑学基本概念的建立（在 1930 至 1940 年间完成）所带来的简单性与一般性方面的进步；另一方面是有界集概念所取得的重要性，它由科尔莫哥洛夫和冯·诺伊曼于 1935 年引入，其在的对偶理论中的根本作用已由麦基{[}213 a and b{]}和格罗滕迪克{[}138 b and c{]}的工作所揭示。最后尤其是，可以肯定，激励这项研究的主要动力来自把理论应用于分析的新可能性，即巴拿赫理论无能为力的那些领域：这里必须提到由克特、托普利茨及其学生们自 1934 年起在一系列论文中发展的序列空间理论{[}184{]}，范塔皮耶“解析泛函”理论的最近建立，尤其是 L. 施瓦茨的广义函数理论{[}280{]}，现代局部凸空间理论在这里找到了一个无疑还远未穷竭的应用领域。

